\section{Conclusion}
In this paper, we have presented \emph{PRDiT} -- a novel and scalable residual Diffusion Transformer model which efficiently solves the 3D CT Volume generation task via two dedicated branches and a two-stage training strategy, consisting of: a Local Denoiser and a Global Residual Transformer. By decoupling the task of modeling coarse local details and refining high-frequency global residuals, PRDiT effectively reduces computational overhead. Extensive evaluations on LIDC-IDRI and Rad-ChestCT datasets demonstrate that our method achieves superior generative performance compared to state-of-the-art methods such as HA-GAN, 3D LDM and WDM-3D, consistently producing lower 3D FID, MMD, and Wasserstein distances. Our experiments further validate the model's scalability, illustrating consistent performance improvements with increasing transformer depth. These results highlight PRDiT's potential as a powerful and efficient method for synthesizing high-resolution 3D medical volumes, with promising applicability to various clinical and diagnostic scenarios.

